% =============================================================================
% D2 v3 — Nature du capital informationnel — réécriture complète 16 mai 2026
% Remplace tout entre \clearpage avant D2 et \clearpage avant D3 (Travail)
% =============================================================================

\subsection{Nature du capital informationnel}\label{sec:d2-capital}

La dimension précédente a examiné les propriétés du bien informationnel et les effets économiques qui en découlent. Le bien est l'objet qui circule~; le capital est le stock qui s'accumule~: les machines, les réseaux, les logiciels, les bases de données, les infrastructures dans lesquels on investit pour produire, stocker ou transmettre de l'information. La question de cette dimension est de savoir si les propriétés du bien suffisent à expliquer les effets économiques du capital, ou si la manière dont le capital est \emph{organisé}~--- sa composition, les conditions d'accès qu'il impose, le degré auquel il est lié à d'autres composants~--- joue un rôle propre. L'enjeu n'est pas abstrait~: les dépenses en calcul informatique liées à l'intelligence artificielle aux États-Unis sont passées de 37~milliards de dollars en~2023 à 219~milliards en~2025, soit une croissance de 144~\% par an, tandis que la production ajustée de la qualité croissait de plus de 2\,000~\% par an sur la même période \parencite{korinek_mckelvey_2026}. L'écart entre ces deux chiffres~--- un facteur seize~--- est exactement ce que la mesure du capital par son prix ne voit pas.

Un cas historique rend la question concrète. La calculatrice mécanique et la tabulatrice à cartes perforées sont deux machines de calcul de la même époque, vendues aux mêmes types de clients~--- les grandes organisations~--- avec une technologie comparable. Pourtant, la première produit un marché fragmenté d'une centaine de firmes en~1904 \parencite{cortada_before_1993}, et la seconde un monopole à 85,7~\% en~1935. La technologie ne diffère pas. Les propriétés du bien~--- analysées dans la dimension précédente~--- ne diffèrent pas non plus. Ce qui diffère, c'est la manière dont le capital est construit autour du bien~: la calculatrice est vendue, le client possède la machine et peut changer de fournisseur~; la tabulatrice est louée, avec un flux de consommables exclusifs et des données accumulées dans un format propriétaire.


\subsubsection{Le champ~: pourquoi la forme du capital est une question}

Le contraste calculatrice/tabulatrice pose un problème que la théorie du capital ne résout pas entièrement.

\markNEO
Une première explication serait que la tabulatrice coûte plus cher ou produit davantage, et que la concentration reflète un avantage de productivité. C'est la logique de la comptabilité de la croissance (\emph{growth accounting})~: le capital contribue à la production en proportion de son coût d'usage \parencite{solow_contribution_1956,jorgenson_retrospective_2008}. \textcite{oliner_resurgence_2000,oliner_explaining_2007} ont montré, sur données américaines, que l'investissement en capital informatique expliquait l'essentiel de l'accélération de la productivité du travail entre 1995 et~2004. Mais cette explication ne fonctionne pas ici~: les coûts fixes des deux machines sont comparables, et rien n'indique que la tabulatrice soit intrinsèquement plus productive. Ce qui diffère n'est pas combien le capital coûte ou combien il produit~--- c'est la manière dont la firme y donne accès\footnote{\textcite{greenwood_long_1997} formalisent la logique du changement technique incorporé dans l'investissement (\emph{investment-specific technical change}). \textcite{falck_digital_2024} héritent de cette approche dans un modèle multi-sectoriel. \textcite{jones_ai_rd_2025} en montre la limite~: même une intelligence machine infinie ne fait que doubler le rythme de progrès quand la complémentarité entre composants est forte~--- c'est la structure du capital, pas sa productivité, qui détermine le plafond.}.

Une deuxième explication serait que la tabulatrice a des propriétés économiques différentes~: plus extensible à grande échelle (\emph{scalable}), plus difficilement récupérable (\emph{sunk}), avec davantage d'externalités. C'est la grille de \textcite{haskel_capitalism_2018}, qui identifient quatre propriétés du capital immatériel. Mais appliquées au contraste calculatrice/tabulatrice, ces propriétés ne discriminent pas~: les deux machines ont une extensibilité comparable, des externalités négligeables, et des complémentarités limitées. L'irrécupérabilité diffère~--- la tabulatrice est plus difficile à abandonner~--- mais Haskel et Westlake la traitent comme une propriété générale de l'immatériel, sans distinguer l'actif isolé (dont l'irrécupérabilité est bornée) du système intégré (dont l'irrécupérabilité est amplifiée par les liaisons entre composants). La grille ne voit pas ce qui fait la différence~: ce ne sont pas les propriétés du bien qui changent, c'est la manière dont le capital est construit autour du bien.

Une troisième explication serait que la tabulatrice crée un verrouillage~: les données accumulées, les compétences formées et les contrats signés rendent le changement trop coûteux. \textcite{hughes_networks_1983} théorise l'inertie (\emph{momentum}) des grands systèmes techniques~; \textcite{unruh_understanding_2000} l'étend au verrouillage énergétique fossile (\emph{carbon lock-in}). Cette explication capture une partie du phénomène, mais pas sa cause~: la calculatrice ne produit aucun verrouillage, la tabulatrice en produit un massif. Les deux pourraient en principe créer un système~; la différence est dans le choix de louer plutôt que vendre, d'imposer un consommable exclusif, d'accumuler les données du client dans un format fermé. Le verrouillage n'est pas un accident de l'histoire~--- c'est un résultat de l'architecture du capital.

\markREG
Le problème de l'agrégation du capital prend ici une forme concrète. \textcite{robinson_production_1953} montrait que la valeur du capital dépend du taux de profit qu'elle est censée déterminer~--- une circularité que le capital physique, suffisamment homogène, permettait de contourner en pratique. Le capital informationnel rend le contournement moins tenable~: un câble sous-marin, un modèle de fondation et un brevet logiciel n'ont pas assez de propriétés communes pour qu'un indice de prix les agrège sans perte. Le contraste calculatrice/tabulatrice illustre le problème au niveau micro~: deux capitaux de même prix et de même productivité produisent des résultats opposés. Le problème a aussi une dimension pratique~: l'activité liée à l'intelligence artificielle est aujourd'hui dispersée dans des dizaines de catégories sectorielles~--- centres de données, édition de logiciels, services professionnels, informatique en nuage~--- ce qui rend le capital informationnel invisible dans les comptes nationaux avant même qu'on tente de l'agréger \parencite{korinek_mckelvey_2026}.

Ce que le prix, les propriétés du bien, et l'inertie systémique seule ne captent pas, c'est ce que l'on appellera ici la \emph{forme} du capital~: sa composition (de quoi est fait le stock~?), son mode d'appropriation (comment la firme contrôle-t-elle l'accès~?), et son degré d'intégration systémique (le capital fonctionne-t-il seul ou dans un système de composants complémentaires~?). Le tableau~\ref{tab:d2-capital} présente ces trois caractéristiques, avec la réversibilité comme propriété résultante, appliquées à six configurations historiques et contemporaines.

\begin{table}[!ht]
\caption{Caractéristiques du capital informationnel~: variation entre configurations.}\label{tab:d2-capital}
\centering\footnotesize
\renewcommand{\arraystretch}{1.35}
\begin{tabularx}{\textwidth}{@{}p{1.6cm}p{2.8cm}p{1.6cm}p{1.6cm}p{1.6cm}p{1.6cm}p{1.6cm}p{1.6cm}@{}}
\toprule
 & \textbf{Question} & \textbf{Calculat.} & \textbf{Tabulat. IBM} & \textbf{Télégr.} & \textbf{Babbage} & \textbf{Data center} & \textbf{SaaS/Cloud} \\
\midrule
Composition & De quoi est fait le stock~? & Bien d'équip. & Machine + cartes + données & Réseau linéaire & Prototype unique & Bâtiment + GPU + modèles + PPA & Logiciel hébergé + données client \\
Mode d'accès & Comment y accède-t-on~? & Vente & Location + consom. exclusif & Réseau~; accès comme service & Subvention & Location (cloud) + API fermées & Abonnement + données hébergées \\
Intégration & Seul ou en système~? & Isolé & Intégré (format propr.) & Réseau (valeur = couverture) & Isolé & 4~couches croisées & Services managés propr. \\
Réversibilité & Peut-on récupérer~? & Oui & Non (coût = re-perforage) & Non (infra. sans alt.) & Non (coût pur) & Non (12~mois à 25~ans) & Non (données captives) \\
\bottomrule
\end{tabularx}
\renewcommand{\arraystretch}{1.0}
\end{table}


\subsubsection{Trois observations}

La grille du tableau fait apparaître trois régularités dans le matériau historique et contemporain.

La première concerne la variance de la forme. La calculatrice et la tabulatrice sont deux machines comparables~; ce qui diffère, c'est le mode d'accès. Le résultat est un marché fragmenté pour l'une, un monopole pour l'autre. Le passage du logiciel sous licence perpétuelle au modèle SaaS dans les années~2010 reproduit le même pattern~: la technologie ERP reste comparable, mais le passage de la vente à l'abonnement avec données hébergées concentre le marché~--- les trois premiers fournisseurs (SAP, Oracle, Microsoft) captent plus de la moitié des revenus d'un marché de 73~milliards de dollars en~2025, dont 70~\% en déploiement cloud\footnote{Mordor Intelligence 2026. SAP a étendu le support de ses versions installées localement à~2033 pour amortir la transition.}. Le cas inverse est tout aussi révélateur~: quand le mode d'accès est ouvert~--- logiciel libre, format standard, données portables~--- le verrouillage ne se forme pas. Linux, PostgreSQL et Apache produisent des marchés fragmentés et contestables \parencite{benkler_wealth_2006,hess_understanding_2007}.

La deuxième concerne la sensibilité à la demande. Le télégraphe sert une demande de coordination (commerce transatlantique, gestion ferroviaire) et produit un capital de réseau qui se substitue au capital physique. La calculatrice sert une demande de traitement interne (assurances, comptabilité) et produit un capital d'outil qui complète le travail. La \emph{Difference Engine} de Babbage ne sert aucune demande effective et produit un coût irrécupérable pur. Le même type de capital a des effets radicalement différents selon la demande qu'il sert. Mais cette observation soulève une objection~: est-ce la demande qui détermine l'effet, ou la technologie~? Le télégraphe est un réseau, la calculatrice est un outil~--- ce n'est pas la même technologie. L'observation ne permet pas d'isoler le facteur demande du facteur technologie. Ce qu'elle établit, c'est que le capital informationnel n'est pas un facteur homogène~--- pas que la demande est le seul facteur discriminant.

La troisième concerne la variance de l'irréversibilité. La calculatrice est substituable~: le client peut changer de fournisseur sans coût significatif. Le centre de données contemporain ne l'est pas~: quatre couches à temporalités incompatibles (bâtiment 25~ans, processeurs graphiques 3 à 5~ans, modèle 12 à 18~mois, contrat d'énergie 10 à 15~ans) sont physiquement inséparables dans l'exploitation. Entre ces deux pôles, l'écosystème CUDA de Nvidia est rendu non substituable par 19~ans d'accumulation logicielle (quatre millions de développeurs, 86~\% du marché des centres de données en~2026\footnote{GPUnex 2026. Aucun hyperscaler n'a réduit significativement ses achats Nvidia.}). Le même gradient s'observe dans le capital industriel numérisé~: un automate programmable isolé est substituable~; intégré dans un système couplant MES, ERP, capteurs et protocoles de communication à des temporalités comparablement dispersées, il ne l'est plus. L'observation est que l'irréversibilité varie continûment entre les configurations, mais les cas disponibles ne permettent pas encore d'isoler le facteur qui en détermine l'amplitude~--- intégration systémique, spécificité physique de l'actif, ou accumulation logicielle.


\subsubsection{Trois conjectures}


\paragraph{La substitution conditionnée par la demande}

\markREG
\textcite{field_magnetic_1992}, dans un article du \emph{Journal of Economic History}, fournit la mesure la plus précise d'une substitution capital informationnel / capital physique. Il montre que le télégraphe a permis aux compagnies ferroviaires américaines de maintenir des voies uniques là où la sécurité aurait exigé un doublement~: le capital informationnel se substitue au capital physique avec un ratio de levier de 1:7 à 1:13\footnote{Le ratio repose sur la comparaison entre la capitalisation de Western Union en~1893 (123~millions de dollars) et le coût estimé du doublement ferroviaire (957~millions). C'est un ordre de grandeur, pas une mesure précise.}. Ce résultat n'entre ni dans la comptabilité de la croissance de \textcite{jorgenson_information_2001}, qui traite le capital informatique comme un facteur homogène, ni dans le cadre tâches-compétences de \textcite{autor_skill_2003}, qui traite la substitution au niveau du poste de travail. Ce que Field mesure est un canal distinct, qui opère au niveau de l'infrastructure.

L'analyse historique suggère que ce canal n'est pas universel. La calculatrice ne substitue rien~: \textcite{reinstaller_historical_2001} documente une corrélation positive de 0,921 entre le stock de capital de bureau et l'emploi de bureau aux États-Unis entre 1870 et~1940~--- le capital réduit le coût unitaire du calcul, la demande de calcul augmente, et l'emploi augmente avec elle. Babbage est à l'autre extrémité~: un investissement sans demande, un coût irrécupérable pur.

L'analyse contemporaine suggère la même conditionnalité. Le mobile bancaire en Afrique subsaharienne (M-Pesa, lancé au Kenya en~2007) sert une demande d'inclusion financière~: le capital informationnel se substitue à l'infrastructure bancaire physique, avec un capital léger et réversible. L'entraînement de modèles de fondation, dont le coût unitaire dépasse cent millions de dollars \parencite{cottier_rising_2025}, sert une demande de capacité cognitive générique, indéterminée dans son usage~--- plus proche de Babbage que de Field.

\paragraph{Observation stylisée} \emph{L'analyse historique et contemporaine suggère que la fonction économique du capital informationnel~--- substitution, complémentarité, coût irrécupérable~--- est conditionnée par la structure de la demande qu'il sert, pas par ses propriétés techniques intrinsèques. Le matériau historique fait apparaître une tripartition~--- coordination, traitement interne, absence de demande effective~--- mais cette tripartition ne couvre pas l'ensemble des configurations contemporaines~: le capital de réseau sert plusieurs types de demande simultanément, le logiciel embarqué sert une demande de performance fonctionnelle d'un bien physique, et la demande de semi-conducteurs avancés est partiellement déterminée par des contraintes géopolitiques. Le mécanisme de conditionnalité par la demande est observable~; sa formalisation complète reste ouverte.} Cette conjecture a un domaine de validité~: elle discrimine tant que le capital informatique est spécifique à certaines demandes. \textcite{restrepo_wont_be_missed_2025} formalise un régime dans lequel le capital computationnel automatise les tâches de goulot d'étranglement elles-mêmes~: les salaires convergent vers le coût computationnel de réplication du travail humain, et la part du travail tend vers zéro. Dans ce régime, la conditionnalité par la demande cesse de discriminer~--- la question se déplace vers la capture des rentes et l'intensité des contraintes physiques.


\paragraph{Le mode d'appropriation comme architecture de marché}

\markINST
Le contraste calculatrice/tabulatrice pose un fait que la littérature ne résout pas entièrement~: ce n'est pas le coût fixe qui discrimine la structure de marché~--- les deux technologies ont des coûts fixes comparables~--- c'est la manière dont la firme organise l'accès au capital. Dans le cas contemporain des modèles de fondation, les deux canaux opèrent simultanément~: le coût fixe est considérable ($>$100~millions de dollars) et crée des barrières indépendamment du mode d'accès \parencite{athey_ai_competition_2025}. La contribution propre de cette analyse est d'identifier le canal du mode d'appropriation comme \emph{distinct} de celui du coût fixe, pas de s'y substituer.

Plusieurs auteurs approchent cette idée sans la formuler sous cette forme. \textcite{de_ridder_2024} formalise la concentration par les coûts fixes des immatériels. \textcite{williamson_economic_1985} traite la spécificité des actifs comme une propriété donnée, pas comme un objet de choix. \textcite{shapiro_varian_1999} analysent les coûts de changement dans un cadre statique de tarification. \textcite{lessig_code_1999} montre que le code est une forme de régulation, mais ne relie pas cette idée à la théorie du capital. Ce que l'on propose d'appeler ici \emph{régime d'appropriation du capital}~--- l'ensemble des modalités par lesquelles l'accès au capital est organisé~--- n'est formalisé sous cette forme par aucun de ces auteurs. Le terme est une proposition de cette thèse, utilisé sous cette réserve. Le régime d'appropriation peut résulter d'un choix stratégique de la firme productrice (IBM louant ses tabulatrices, les hyperscalers tarifant l'extraction de données), d'une contrainte technologique (la captivité des designs dans un n\oe ud de fabrication chez TSMC), d'une configuration de chaîne de valeur (le verrouillage créé par l'intégrateur industriel, non par le fabricant du composant), ou d'un cadre réglementaire (les licences spectrales et les obligations de dégroupage dans les télécommunications). La conjecture porte sur l'\emph{effet} du régime d'appropriation sur la structure de marché, quel que soit l'agent ou le mécanisme qui le détermine.

Une raison théorique de l'importance de ce design tient à ce que \textcite{haskel_capitalism_2018} appellent, sans le formaliser, la contestabilité de l'immatériel~: les externalités brouillent la propriété, les complémentarités rendent chaque composant indissociable du reste, et les protections juridiques ne suffisent pas à garantir le contrôle\footnote{\textcite{dumez_croissance_2018} souligne ce point dans sa recension de Haskel et Westlake.}. Le cas IBM l'illustre~: les brevets ont été contestés avec succès au tribunal (\emph{United States v.~IBM}, 1935), mais le verrouillage a survécu au jugement~--- les coûts de migration étaient incorporés dans les données accumulées sur un format fermé, pas dans le brevet.

Le passage au SaaS et le cas du cloud public confirment le même pattern dans le registre contemporain. Les trois hyperscalers (Amazon Web Services, Microsoft Azure, Google Cloud Platform) contrôlent 67 à 68~\% du marché mondial de l'infrastructure cloud, qui atteint 395~milliards de dollars en~2025\footnote{Synergy Research Group, quatrième trimestre~2025. AWS~: environ 31~\%, Azure~: 24~\%, GCP~: 12~\%.}. La structure tarifaire est révélatrice~: transférer des données vers le cloud est gratuit~; les en sortir coûte 0,01 à 0,12~dollar par gigaoctet. L'asymétrie reproduit la structure du modèle IBM~--- gratuit à l'entrée, coûteux à la sortie~--- transposée à l'échelle planétaire.

Le passage au modèle cloud a une conséquence comptable qui amplifie le verrouillage. Pour le client, le capital informatique cesse d'apparaître au bilan~: ce qui était un investissement capitalisé (serveurs, licences perpétuelles) devient une dépense courante d'abonnement. Le capital est toujours là~--- il est dans le bilan de l'hyperscaler~--- mais il est invisible dans celui du client. Cette transformation CAPEX-OPEX rend la dépendance plus difficile à percevoir et à évaluer~: un directeur financier voit une ligne de charges récurrentes, pas un stock d'actifs captifs dont le coût d'extraction excède le coût de migration. Le verrouillage n'est pas seulement technique (données hébergées, API fermées)~--- il est aussi \emph{comptable}, au sens où le traitement en charges masque l'accumulation d'un actif non portable chez le fournisseur.

Le cas négatif est tout aussi important. Quand le mode d'accès est ouvert~--- logiciel libre, format standard, données portables~--- la condition d'activation n'est pas remplie~: le client n'accumule pas d'actif captif, et le verrouillage ne se forme pas. Linux, PostgreSQL et Apache produisent des marchés fragmentés et contestables \parencite{benkler_wealth_2006}. \textcite{hess_understanding_2007} formalisent ce régime comme une gouvernance par les communs, distincte du marché et de l'État.

\paragraph{Conjecture} \emph{L'analyse suggère que le mode d'appropriation du capital informationnel est un déterminant endogène de la structure de marché, distinct du coût fixe et de la spécificité des actifs. La condition d'activation est que le client accumule un actif non portable~--- données, configurations, intégrations~--- dont le coût d'extraction excède le coût de migration. Quand cette condition est remplie, le mode d'appropriation concentre le marché~; quand le mode d'accès est ouvert, elle ne se crée pas.} La vérification systématique exigerait un panel de cas variant le mode d'accès à technologie constante et la mesure de l'effet sur l'indice de concentration. \textcite{korinek_lockwood_2025} en tirent une conséquence fiscale~: quand la part du travail décline, les rentes économiques deviennent la cible prioritaire, et le design du capital détermine qui les capture. \textcite{korinek_mckelvey_2026} documentent un phénomène qui renforce ce mécanisme~: le déflateur d'inférence, ajusté de la qualité, baisse de 94~\% par an~--- chaque dollar de calcul produit seize fois plus qu'un an auparavant. Les termes de l'échange du secteur se dégradent à mesure que la production explose mais que les prix s'effondrent. La structure est analogue à ce que \textcite{bhagwati_immiserizing_1958} appelle une \og~croissance appauvrissante~\fg{} dans le commerce international~: un pays exportateur dont la production croît mais dont les termes de l'échange se dégradent plus vite. L'emprunt est métaphorique~--- le producteur d'inférence n'est pas un pays, et le mécanisme de transmission diffère~--- mais la structure logique est commune~: la valeur migre vers l'aval (l'usage) et la capture de rentes en amont exige un contrôle de l'accès, pas de la production. Dans ce contexte, le régime d'appropriation du capital n'est pas un choix stratégique parmi d'autres~: c'est le mécanisme principal par lequel la firme capture de la valeur quand le prix unitaire de sa production tend vers zéro. IBM capturait la valeur par les cartes perforées quand le coût de la machine devenait secondaire~; les hyperscalers capturent la valeur par les API fermées et les données hébergées quand le coût d'inférence tend vers zéro.


\paragraph{L'irréversibilité comme gradient}

\markNEO \markSTS
La tradition du capital incorporé (\emph{embodied technical change}), depuis \textcite{solow_investment_1960} jusqu'à \textcite{johansen_substitution_1959}, établit que le progrès technique est matérialisé dans les nouvelles générations de machines~: une fois installé, le capital incorpore la technologie de sa date de naissance et ne s'améliore plus. C'est le fondement du modèle où le capital est flexible à l'installation et rigide ensuite (\emph{putty-clay}). Ce que le capital informationnel intégré incorpore n'est pas seulement une technologie~--- c'est une architecture d'accès, un écosystème logiciel, des formats de données, des contrats et des compétences dont la rigidité excède celle de la machine individuelle. L'entraînement des modèles d'intelligence artificielle produit une catégorie de capital que \textcite{korinek_mckelvey_2026} appellent \emph{capital de modèle} (\emph{model capital})~--- un actif immatériel qui incorpore la technologie de sa date d'entraînement et ne s'améliore plus ensuite, conformément à la logique du capital incorporé. Mais sa vitesse d'obsolescence est sans précédent~: le coût de calcul nécessaire pour atteindre un niveau de performance donné baisse d'environ deux tiers par an \parencite{ho_algorithmic_2024}, et le déflateur d'inférence ajusté de la qualité baisse de 94~\% par an~--- bien au-delà des 30 à 50~\% par an documentés pour les semi-conducteurs par \textcite{byrne_semiconductor_2018}. Le capital de modèle n'est ni du tangible (pas de machine physique) ni de l'immatériel classique au sens de Haskel-Westlake (pas un brevet à durée de vie longue, pas du logiciel au sens comptable). C'est un objet économique nouveau dont la vitesse de transformation défie les catégories existantes.

Cette vitesse crée un paradoxe. Le capital informationnel intégré est simultanément hyper-réversible dans ses couches rapides (un modèle est obsolète en 12 à 18~mois) et hyper-irréversible dans ses couches lentes (le bâtiment dure 25~ans, le contrat d'énergie 15~ans). Et le paradoxe se renforce~: plus les couches rapides se transforment vite, plus les couches lentes \emph{contraignent}~--- parce qu'on ne peut pas recâbler, relocaliser ou reconvertir l'enveloppe physique au rythme de l'obsolescence logicielle. Le gradient isolé/intégré n'est pas seulement un gradient de structure~--- c'est un gradient de \emph{tension} entre vitesses incompatibles, et cette tension s'aggrave à mesure que la vitesse d'obsolescence des couches rapides augmente.

L'irréversibilité de l'investissement crée une valeur d'option de l'attente \parencite{dixit_pindyck_1994}~: le seuil d'investissement optimal est, pour des paramètres typiques, le double de ce que la valeur actualisée nette prédirait. Ce cadre traite un actif isolé. À l'autre pôle, \textcite{hughes_networks_1983} théorise l'inertie des grands systèmes techniques~; \textcite{unruh_understanding_2000} l'étend au verrouillage énergétique fossile. L'irréversibilité n'est plus seulement financière~--- elle est systémique.

Le gradient va de l'actif isolé (la calculatrice~: réversible, substituable) au système intégré dans un complexe techno-institutionnel (le chemin de fer coordonné par McCallum~: irréversible, co-évolutif avec les institutions). L'écosystème CUDA de Nvidia occupe une position intermédiaire révélatrice~: le matériel est rendu non substituable par l'écosystème logiciel accumulé en dix-neuf ans, et les coûts de migration se cumulent de manière multiplicative. Le centre de données hyperscale est au pôle intégré~: quatre couches à temporalités incompatibles, physiquement inséparables.

Ce gradient n'est pas celui de l'irrécupérabilité au sens de \textcite{haskel_capitalism_2018}~: Haskel et Westlake traitent l'irrécupérabilité comme propriété de l'immatériel en général, sans distinguer l'actif isolé du système intégré. Le gradient proposé ici est isolé/intégré, pas matériel/immatériel~--- c'est une proposition de la thèse, pas un résultat établi. Ce gradient a un analogue structurel dans le paramètre~$\theta$ de \textcite{jones_ai_rd_2025}~: plus le capital est intégré dans un système de composants complémentaires non automatisables, plus le plafond de productivité est bas. L'objet n'est pas le même~--- $\theta$ mesure la complémentarité entre tâches, le gradient mesure la complémentarité entre couches~--- mais la structure logique est commune.

Les contre-arguments sont à prendre au sérieux. \textcite{greenstein_lock_in_1997} montre que le verrouillage informatique a historiquement procédé de la compatibilité logicielle, non de l'irréversibilité physique~--- et qu'il a été plus léger que prévu. \textcite{liebowitz_margolis_1995} contestent que la dépendance de sentier forte soit empiriquement fréquente. La question ouverte est de savoir si la matérialisation croissante de l'infrastructure~--- centres de données à l'échelle du gigawatt, contrats énergétiques de quinze ans~--- déplace le verrouillage computationnel du registre Greenstein (architectural, réversible à un coût) vers le registre Hughes-Unruh (systémique, structurellement irréversible). Une note récente du FMI souligne un canal complémentaire~: les anticipations de croissance liées à l'IA affectent les taux réels et la dynamique de dette \emph{avant} que les gains de productivité se matérialisent \parencite{barhoumi_imf_ai_2026}, créant un verrouillage financier qui amplifie le verrouillage technique.

Une interaction entre conjectures mérite d'être signalée. Le mode d'appropriation peut amplifier l'irréversibilité~: un capital loué avec données hébergées est plus irréversible qu'un capital vendu avec données portables, indépendamment de son degré d'intégration technique. Les deux conjectures interagissent dans le cas du cloud, où le mode d'accès (API fermées, données hébergées) renforce l'irréversibilité technique (couches à temporalités croisées).

\paragraph{Conjecture} \emph{L'analyse suggère que l'irréversibilité du capital informationnel varie sur un gradient continu, déterminé par le degré d'intégration systémique~--- complémentarité entre couches, liaison physique, co-évolution institutionnelle~--- pas par la nature matérielle ou immatérielle de l'actif.} Les conditions d'activation sont~: forte complémentarité technique entre couches, rythmes de dépréciation hétérogènes, et co-évolution avec des institutions et des contrats. Quand l'une au moins fait défaut, le gradient ne s'active pas. Les contre-arguments de Greenstein et Liebowitz-Margolis imposent de ne pas présumer que le gradient s'amplifie~; la question reste ouverte.

Le gradient proposé capture l'irréversibilité \emph{systémique}~--- celle qui naît du couplage entre couches complémentaires à temporalités hétérogènes. Il ne capture pas l'irréversibilité de \emph{spécialisation}~--- celle qui naît de la calibration physique d'un actif pour un usage unique, indépendamment de toute intégration dans un système plus large. Un site de fabrication de semi-conducteurs avancés (coût unitaire supérieur à vingt milliards de dollars) est irréversible non par couplage multi-couches mais par spécificité physico-chimique~: les salles blanches, les processus de lithographie et les équipements sont dédiés à un n\oe ud technologique précis et ne sont pas reconvertibles. Les deux types d'irréversibilité coexistent dans le capital informationnel lourd~; la conjecture ici formulée ne porte que sur le premier.


\subsubsection{Fait stylisé}

La forme du capital informationnel~--- sa composition, son mode d'appropriation, son degré d'intégration systémique~--- est conditionnée par la structure de la demande qu'il sert et par le contexte institutionnel dans lequel il s'inscrit. L'analyse historique et contemporaine suggère que cette forme conditionne à son tour le type d'effet économique (substitution, complémentarité, verrouillage), le degré d'irréversibilité, et la structure de marché qui en résulte. Quand la technologie est constante (calculatrice vs tabulatrice, licence vs SaaS), c'est la forme qui discrimine. Quand la forme est constante, c'est la demande qui discrimine (télégraphe vs calculatrice). Dans le cas général, les deux interagissent.

Le gradient va de l'outil isolé (réversible, marché fragmenté) au système intégré (irréversible, concentrant). Ce gradient opère de 1820 au centre de données contemporain, et discrimine à chaque époque. Les segments les plus récents de l'infrastructure numérique~--- là où le capital est le plus intégré systémiquement~--- poussent cette logique jusqu'au point où le capital informationnel acquiert les propriétés du capital d'infrastructure lourde. C'est un cas extrême du gradient, pas un régime universel~: le mobile en Afrique subsaharienne, le logiciel libre et l'IA applicative montrent que d'autres formes, plus légères et plus réversibles, coexistent.

Ces conjectures ont un domaine de validité explicite. La conditionnalité par la demande discrimine tant que le capital informatique est spécifique~; dans le régime d'automatisation généralisée formalisé par \textcite{restrepo_wont_be_missed_2025}, elle cesse de discriminer et la question se déplace vers la capture des rentes et l'intensité des contraintes physiques~--- énergie, terre, matériaux~--- formalisées par le paramètre~$\theta$ de \textcite{jones_ai_rd_2025}. Les conjectures sur le mode d'appropriation et sur l'irréversibilité restent opérantes dans les deux régimes. Elles atteignent cependant une limite plus fondamentale si le capital informationnel cesse d'être un instrument de la production humaine pour devenir un système productif autonome. \textcite{korinek_mckelvey_2026} proposent de traiter l'économie de l'IA comme une économie séparée~--- une \og~île~\fg{} qui importe de l'électricité, des puces et du travail, et exporte de l'inférence. Dans ce cadre, le capital informationnel n'est plus un stock dans le bilan d'une firme~--- c'est l'interface entre deux systèmes productifs fonctionnant à des vitesses radicalement différentes (2~\% de croissance côté humain, 2\,600~\% côté IA). Si cette transition se matérialise, la forme du capital n'est plus seulement une propriété du stock~--- c'est l'architecture des termes de l'échange entre production humaine et production computationnelle. Les conjectures ici formulées opèrent en deçà de ce seuil~; au-delà, la catégorie même de capital informationnel devrait être repensée. Une rupture technologique qui découplerait radicalement les couches~--- par exemple un calcul neuromorphique entièrement embarqué, sans infrastructure centralisée~--- en invaliderait également le diagnostic.
